\chapter{Four banking translations before formal causal analysis}
\label{exec13:l2-translation}

A request such as ``Should we improve retention?'' is not yet a causal question. The action may be a rate offer, a service contact, a product change or a budgeted policy. The unit may be an account, customer, household or branch. The outcome may be balance retention, funding cost, repayment or relationship value. An assistant should expose these alternatives and ask for the missing decision; a formal engine cannot repair a target chosen incorrectly before the graph is built. The four translations below keep the business motivation while making different estimands visible. Their formal discipline follows the separation of a target trial, its eligible population and its observed emulation in \citet{hernan2020}; the banking choices themselves are proposals, not observations.

\section{Deposit pricing: acquisition, retention or net funding cost}

Suppose a treasury manager considers a 25-basis-point deposit-rate increase. Translation A asks for the effect of offering the rate to every eligible account on thirty-day average balance. Translation B asks for the effect on balances among accounts that historically received similar offers. Translation C asks for net funding cost, including paid interest, incentive expense, repricing of existing balances and any value assigned to retained relationships. Translation D asks for a fixed-budget allocation policy that chooses which accounts receive the offer. These are four populations or outcomes, not four estimators of one effect.

Let $U$ be the account, $P$ the accounts eligible under the written rule, $A$ the offer, and $Y^B(a)=H^{-1}\sum_{d=1}^H B_d(a)$ the average daily balance over $H=30$ days, measured in dollars. The corresponding balance-time outcome is $S(a)=H Y^B(a)$, measured in dollar-days. The all-eligible balance target is $\mathbb E_P[Y^B(1)-Y^B(0)]$. A monetary value target instead specifies a valuation $M(a)$ in dollars and compares $\mathbb E_P[M(1)-M(0)-\{C(1)-C(0)\}]$, where $C(a)$ is monetary cost outside $M(a)$. Costs already inside $M(a)$ must not be subtracted again. If $Y^C(a)$ already contains net funding cost, the objective is to minimize $\mathbb E_P[Y^C(a)]$ and a negative $\mathbb E_P[Y^C(1)-Y^C(0)]$ favors the offer. The policy target additionally requires a feasible policy class and an explicitly costed comparator. A realized post-offer balance may enter the monetary outcome but is not automatically a pre-offer policy weight. A customer identifier does not establish economic ownership when accounts are shared or linked. Section~\ref{exec17:l2-monetary} derives a dimensionally consistent valuation and its alternatives.

For example, let $W$ be the frozen pre-offer balance and segment record, $e(W)=\Pr(A=1\mid W)$, and $m_a(W)=\mathbb E(Y^B\mid A=a,W)$. If consistency, conditional exchangeability, positivity and the declared no-interference boundary hold, the all-eligible target has the identifying expression
\[
 \psi_P=\int\{m_1(w)-m_0(w)\}\,dP(w).
\]
This equality is conditional on the assumptions; a fitted model cannot verify them. If offers are absent for a high-balance stratum, restricting the integral to $0<e(W)<1$ changes its averaging measure. Likewise, a policy $\pi(W)$ under a contact budget $B$ asks for $\mathbb E_P\{Y^C(\pi(W))-Y^C(\pi_0(W))\}$ subject to $\mathbb E_P[C(\pi(W))]\leq B$. The target needs a specified comparator $\pi_0$, a cost convention and an availability clock before any optimization result has a business meaning. The target-versus-estimator distinction also appears in the inspected formulation literature \citep{exec13_formulation_targets}.

A valid clock records when rate eligibility, customer status and balance information were available to the offer system. A statement generated after the decision cannot be used as if it were known before it. Diagnostics include overlap by eligibility stratum, repricing spillovers, concurrent campaigns and whether the offer itself changes future eligibility. A narrower overlap target can be useful, but the report must display the excluded accounts and say that the original all-eligible target remains unresolved.

Consider two equally common eligible strata with conditional effects of $2$ and $8$ balance units. Their all-eligible effect is $5$. If the offer probabilities are $0.1$ and $0.5$, overlap weights $e(1-e)$ are $0.09$ and $0.25$; their normalized average is $(0.09\cdot2+0.25\cdot8)/(0.09+0.25)=109/17\approx6.41$, not $5$. This is exact authored arithmetic, illustrating why a desirable stable-overlap answer cannot be substituted for the requested all-eligible effect. The example asserts no bank offer probabilities.

\section{Arrears servicing: offer, acceptance or dynamic policy}

A collections manager asks whether an early restructuring review should be offered to an account entering arrears. Translation A targets an offer at the first qualifying arrears date. Translation B targets acceptance among customers who take up the review. Translation C targets the effect of a fixed sequence of contacts. Translation D targets a dynamic policy that uses evolving hardship and payment history. Acceptance is downstream of the offer; it is not the same intervention. A static binary ATE cannot silently answer a longitudinal policy question.

Let $X_t$ contain information available at the decision time, $A_t$ the offer, $L_{t+1}$ later hardship information and $Y$ discounted net recovery over 180 days. A total-effect question should not condition on a post-offer descendant $L_{t+1}$ as an ordinary baseline covariate. A policy question needs the action history, availability of each $L_t$, censoring and a comparator policy. If several accounts share a household or collection team, account rows are not automatically independent units. A field entered on day 11 cannot govern a day-10 action even if its event date is day 9.

For a two-decision example let $H_0=X_0$ and $H_1=(X_0,A_0,L_1)$, with a regime $g=(g_0,g_1)$ where $A_0=g_0(H_0)$ and $A_1=g_1(H_1)$. Its target is $V(g)=\mathbb E[Y(g)]$ on the declared first-arrears cohort. Sequential exchangeability, consistency, positivity along the regime and a justified censoring law are needed for a longitudinal identification formula. For the discrete case, write $h=(x,g_0(x),l)$, $p(x)=\Pr(X_0=x)$ and
\begin{align*}
 m_g(h)&=\mathbb E[Y\mid H_1=h,A_1=g_1(h)],\\
 p_g(l\mid x)&=\Pr(L_1=l\mid X_0=x,A_0=g_0(x)).
\end{align*}
Then the longitudinal formula is
\[
 V(g)=\sum_{x,l}m_g(x,g_0(x),l)\,p_g(l\mid x)\,p(x).
\]
The formula illustrates the clock: $L_1$ may guide $A_1$, but cannot retroactively guide $A_0$. It also shows why a coefficient from a regression that conditions on $L_1$ cannot be relabeled the total first-offer effect. This is a project specialization of the longitudinal identification argument in \citet{hernan2020}.

Failure cases are substantive. A model that predicts acceptance can be accurate while answering no offer-effect question. A model that trims accounts with poor overlap and continues to report an all-account effect changes the population. A missing policy-history component is an unsupported target, not a zero effect. The remedy is to return a typed alternative or refusal and request the missing policy, clock or population decision.

\section{Credit underwriting: approval threshold or feasible decision rule}

A credit team may ask whether lowering an approval threshold changes loss. Translation A is the effect of offering credit to applicants between the old and new thresholds. Translation B is the effect on repayment conditional on approval. Translation C is the effect on portfolio loss including volume, utilization and capital cost. Translation D is a policy value under a capacity or risk budget. Approval is an action selected by a threshold, while repayment is observed only under selected applicants; conditioning on approved cases changes the target.

The unit is an application at the decision date. The action must specify amount, term and price, not only a binary approval label. Potential outcomes include repayment, loss, utilization and capital usage over a fixed horizon. A threshold change can alter composition, interference through capacity and future collection intensity. The information contract must distinguish bureau vintage, application-time fields, later delinquency and revised labels. Diagnostics include support around the threshold, manipulation or anticipatory behavior, policy feasibility, missing outcomes and whether the threshold is a rule or a human override.

A regression discontinuity or instrumental design requires its own continuity, exclusion and assignment assumptions; a predictive score does not supply them. If the target is an approval-policy value, the report should retain the policy constraints and loss function rather than output a conditional repayment coefficient. A public or generated threshold mechanism can test translation and clock preservation before requesting bank outcome truth.

To see the selection issue, let $D(z)$ denote approval under threshold policy $z$ and let $Y(1)$ denote repayment under a specified offered loan. The mean $\mathbb E[Y(1)\mid D(z)=1]$ describes the applicants selected by $z$; it is not $\mathbb E[Y(1)]$ for all applicants. Changing $z$ can change that conditioning population even if no individual repayment response changes. Near a score cutoff $c$, the discontinuity $\lim_{s\downarrow c}\mathbb E(Y\mid S=s)-\lim_{s\uparrow c}\mathbb E(Y\mid S=s)$ may have a local causal interpretation only under the design's continuity and manipulation restrictions; it is not the portfolio effect of moving the entire threshold. A threshold policy must count declined applications with their opportunity cost and define the loan terms and missing counterfactual repayment law. The local-design distinction is consistent with the audited discontinuity and target-trial sources \citep{hernan2020}.

\section{Branch campaign: donor counterfactual or randomized rollout}

A branch manager asks whether a campaign increased new accounts. Translation A is a single-branch time-series counterfactual using unaffected donors. Translation B is a panel effect across branches with staggered launch dates. Translation C is a cluster-randomized rollout. Translation D is the value of changing campaign intensity under a constrained budget. The outcomes, donor assumptions and uncertainty objects differ.

For a treated branch $Y_t$, donor series $X_t$ and intervention date $T_0$, a structural time-series method predicts $Y_t(0)$ from pre-period history and controls. It requires controls unaffected by the campaign and a stable relationship after $T_0$. Perfect pre-period fit does not establish either restriction. Shared marketing, customer migration or national pricing can contaminate donors. A randomized rollout addresses a different question with a different cost and design; panel methods require their own assignment and trend assumptions.
The inspected structural-time-series and synthetic-control sources use unaffected predictors and design information in their respective comparisons \citep{exec13_ci_brodersen2015,abadie2010}; neither makes a pre-period fit diagnostic a proof of post-period donor invariance.

An exact construction makes the risk concrete. Let $A_t=1$ on four campaign days, $X_t=x_t^0+15A_t$ and $Y_t=100+2x_t^0+20A_t$. The pre-period relationship is $Y=100+2X$. Continuing it post-period gives a contrast of $-10$ per day, although one structural law has a true campaign effect of $20$ and a contaminated donor. A second law produces the same twelve observed rows and a true effect of $-10$ with an unaffected donor. No estimator can distinguish these laws without a restriction or experiment. With a stipulated USD1 value per application and USD20 cost, the four-day net values are +60 and -60. These are authored mechanisms, not bank evidence.

In the first law, the observed post-period contrast is $Y_t-(100+2X_t)=20-2(15)=-10$, while the counterfactual holding the donor's untreated path fixed is $Y_t(0)=100+2x_t^0$ and the true effect is $20$. In the second law, define an unaffected donor $\widetilde x_t^0=X_t$ and an outcome $Y_t=100+2\widetilde x_t^0-10A_t$. Every observed $(A_t,X_t,Y_t)$ row is then identical across the two laws, including the eight pre-period and four post-period rows, yet the effect is $-10$. At USD1 per application and a single USD20 campaign cost, $4(20)-20=60$ and $4(-10)-20=-60$. The sign of the investment decision is therefore not identified from these rows alone. The contrasting laws are reproducible algebra, not a claim that either mechanism occurs at a bank.

\section{The formulation record and formal handoff}

The declaration should contain $Q=(U,P,A,A_0,Y,H,\nu,\ell,I)$: unit, population, action, comparator, outcome, horizon, averaging measure, loss and information clock. It should preserve candidate questions when the business sentence is ambiguous. A schema check verifies fields and types; a semantic check compares them with the brief and dictionary; a formal check applies the supplied graph and method. A passed schema is not causal truth, and a failed backend is not nonidentification.

The order of handoff matters. First, freeze $Q$ and the candidate alternatives with a human owner. Second, attach each proposed graph arrow to a source, temporal ordering and an uncertainty label. Third, ask the formal engine which estimand, if any, follows from the supplied graph and assumptions. Fourth, run estimators only for supported targets and return precision, sensitivity and refusal states. An arrow suggested by retrieval is a hypothesis until its substantive owner accepts it; conversational consistency is not graph accuracy \citep{v12_la_2305_00050,exec13_hallucination}. Formal reasoning benchmarks and narrated synthetic cases help check the translation boundary, but their scores do not validate bank assignment truth \citep{v12_la_2312_04350,v12_la_2607_08093}.

The deterministic baseline is an analyst interview plus fixed arithmetic and graph fixtures. An assisted system can retrieve evidence, draft alternatives, propose a graph and prepare typed calls. It should be scored on target fidelity, clock fidelity, preserved ambiguity, valid abstention, numerical correctness only when the target is identified, and complete failure accounting. All requests belong in the denominator, including refusals and unsupported policies. Public data and generated structural mechanisms are appropriate proposal evidence; authentic bank truth and live provider performance remain later conditional evidence.

The immediate prototype should compare the four translations against the structured baseline using fixed prompts, scenario-level splits, adverse cases and a stopping rule. A later bank study is justified only if the prototype isolates a decision-relevant uncertainty and makes its value of information plausible. No fluent report can substitute for a population decision, assignment record, clock or expert acceptance.
